% =====================================================================
%  introduction.tex — מבוא
% =====================================================================
\chapter*{מבוא: עיר אחת, שני מינים}
\addcontentsline{toc}{chapter}{מבוא}
\markboth{מבוא}{}
\thispagestyle{plain}

השעה \en{06:52}. אתם עדיין בפיג\gr מה, הקפה עוד לא מוכן, ומישהו עומד ליד הדלת ומסתכל עליכם במבט שאומר
"זה הרגע הכי חשוב בחיים שלי". מבחינתו זה באמת כך: הטיול של הבוקר הוא העיתון, הרשת החברתית
וחדר הכושר שלו, הכול ביחד.

החיים עם כלב בעיר הם שותפות בין שני מינים שחווים את אותו רחוב בצורה שונה לגמרי. אתם רואים
מעבר חצייה, אוטובוס ובית קפה. הכלב מריח את ההיסטוריה של כל עמוד חשמל, שומע את האופניים החשמליים
הרבה לפני שאתם שמים לב אליהם, ומבין מצוין שהשקית המרשרשת בכיס שלכם עשויה להכיל משהו טעים.

\secnotoc{איך לקרוא את המדריך}
כל פרק עומד בפני עצמו, אפשר לקפוץ ישר למה שבוער. לאורך הדרך תפגשו תיבות מידע בצבעים שונים,
וכל צבע מסמן סוג אחר של מידע:

\begin{itemize}
  \item \textcolor{sagedark}{\textbf{טיפ חשוב:}} המלצה מעשית שכדאי ליישם כבר בטיול הבא.
  \item \textcolor{navy}{\textbf{כדאי לדעת:}} רקע מקצועי קצר שעוזר להבין \emph{למה}.
  \item \textcolor{coraldark}{\textbf{טעות נפוצה:}} דברים שכולנו עשינו לפחות פעם אחת.
  \item \textcolor{golddark}{\textbf{שימו לב:}} עניינים של בטיחות.
  \item \textcolor{sagedark}{\textbf{טיפ מהשטח:}} הניסיון המצטבר של בעלי כלבים עירוניים.
  \item \textcolor{alert}{\textbf{מידע חירום:}} מצבים שבהם לא מחכים, אלא פונים לווטרינר.
\end{itemize}

\begin{notebox}[הבהרה חשובה]
המדריך נועד להעשרה ולהתמצאות כללית בלבד. הוא לא מאבחן, לא מטפל ולא מחליף ייעוץ אישי אצל
וטרינר או מאלף מוסמך. גם ההמלצות לגבי משכי טיול הן דוגמאות ולא הוראות: גיל, גזע, בריאות ומזג
של כל כלב משנים את התמונה.
\end{notebox}

\secnotoc{מה בפנים}
\textbf{פרק~1} עוסק בראש של הכלב, \textbf{פרק~2} במפה של השכונה, \textbf{פרק~3} בתיק ובשעון,
ו\textbf{פרק~4} ברגעים שבהם משהו משתבש. בסוף המדריך תמצאו מילון מונחים, דף מסכם למקרר,
רשימת מקורות ועמוד שמתאר איך הספר נבנה: במערכת הסדר \en{LaTeX}, בקוד \en{Python} ובשפת השרטוט
\en{TikZ}, ובעזרת שלושה סוכני \en{Claude Code} שעבדו כצוות.

\begin{fieldbox}
הכלב שלכם לא קרא את המדריך הזה. הוא מומחה לעצמו, ולכן כל טיפ כאן שווה בדיוק כמו המידה שבה הוא
מתאים לכלב הספציפי שעומד עכשיו ליד הדלת.
\end{fieldbox}
\clearpage
